\section{Introduction}
\label{sec:introduction}

\subsection{The Problem That Everyone Agrees Is a Problem}

Gerrymandering is the one scandal in American politics that enjoys truly
bipartisan outrage---when the other party is doing it.
Republicans sued over Maryland's congressional map in 2011, calling it a
``drastic partisan gerrymander'' engineered to eliminate a Republican
incumbent.\footnote{%
  \textit{Benisek v.\ Lamone}, 585 U.S. 155 (2019).
}
Democrats sued over North Carolina's map in 2019, securing a ruling from
the state Supreme Court that it violated the state constitution.\footnote{%
  \harper{}, N.C. Supreme Ct., 2022.
}
Republicans sued over New York's map in 2022 and won.\footnote{%
  \harkenrider{}, N.Y. Ct.\ of App., 2022.
}
The litigation never ends because the maps never stop being drawn by people
with a direct interest in the outcome.

The Supreme Court has declined to help.
In \rucho{} v.\ \textit{Common Cause} (2019), the Court held that partisan
gerrymandering presents a political question beyond the federal courts'
reach.\footnote{%
  \rucho{}, 588 U.S. 684, 701 (2019): ``Excessive partisanship in
  districting leads to results that reasonably seem unjust.
  But the fact that such gerrymandering is `unjust' does not mean that it is
  unconstitutional or that it is something federal courts can provide a remedy
  for.''
}
That decision has not ended the problem; it has simply moved the litigation
to state courts, which now manage fifty simultaneous and inconsistent lines
of redistricting jurisprudence.

Meanwhile, the algorithmic solutions proposed by academics and reform advocates
have stalled.
Not because the math is wrong---it is not---but because none of the proposed
algorithms is grounded in a principle that both parties have already accepted.
Every proposed algorithm is a practitioner's choice.
And a practitioner's choice, however mathematically elegant, is a political choice.

This Article proposes a different approach.
It argues that the first principles underlying the existing constitutional
framework---principles that Republicans, Democrats, and the Supreme Court have
all already endorsed without dissent---compel a specific algorithm.
The algorithm is \ar{}.
The principle is \hh{} apportionment applied at the intrastate level.
The law is the \dia{}.

\subsection{Why Algorithmic Solutions Have Stalled}

The academic redistricting literature is rich.
Markov chain Monte Carlo methods \citep{deford2019,metric2018} can generate
large ensembles of valid maps and flag outliers.
Optimization methods \citep{gurnee2021} can minimize specific measures of
unfairness.
Column-generation heuristics can satisfy complex multi-constraint formulations.
Computer scientists have published hundreds of papers on redistricting algorithms
since 2010.

None of this has produced a federal law.
The reason is structural: every proposed algorithm requires a practitioner to
choose which algorithm to use.
That choice is itself a political act.

The empirical evidence is stark.
In Wisconsin, which has eight congressional seats and has been a national
flashpoint for redistricting litigation, the choice of algorithm determines
whether Democrats win four seats or three---an eight-seat delegation swinging
on a methodological choice made before any data is examined.
The bakeoff results documented in \citet{b0paper} show that running the
unweighted bisection algorithm (no geographic weights) on Wisconsin's
2020 census tracts produces a 4~Democrat / 4~Republican split with a partisan
gap of $-0.3$~percentage points.
Running GeoSection---a mathematically sophisticated, geographically rigorous
algorithm developed in \citet{b8paper}---on the same data produces a
3~Democrat / 5~Republican split with a partisan gap of $-12.8$~percentage points.

Neither algorithm has political inputs.
Neither algorithm is trying to advantage one party.
The 12.5-percentage-point difference is entirely an artifact of which algorithm
was chosen.

This is the compactness-proportionality paradox: the choice between
mathematically valid, politically neutral algorithms is itself a political
choice.
The only resolution is to mandate the algorithm that is prescribed by
first principles, not selected by practitioners.
And there is exactly one such algorithm: \hh{} apportionment, which is already
the constitutionally mandated method for the structurally identical
interstate problem.

\subsection{The Argument in Brief}

The argument in this Article has four steps, each resting on a principle
that has already been accepted without meaningful dissent.

\paragraph{Step~1: Equal population.}
\wesberry{} v.\ \textit{Sanders} (1964) held that congressional districts
must have equal population.\footnote{%
  \wesberry{}, 376 U.S. 1, 18 (1964): ``as nearly as is practicable one man's
  vote in a congressional election is to be worth as much as another's.''
}
\karcher{} v.\ \textit{Daggett} (1983) operationalized this as effectively
zero tolerance for population deviations in congressional districts.\footnote{%
  \karcher{}, 462 U.S. 725, 730 (1983): ``the `as nearly as practicable'
  standard requires that the State make a good-faith effort to achieve precise
  mathematical equality.''
}
This is not controversial.
Both parties accept population equality as the floor.

\paragraph{Step~2: Huntington-Hill is the agreed interstate method.}
Article~I, \S2 requires that seats be apportioned among states
``according to their respective Numbers,'' and Congress has operationalized
this via the Huntington-Hill equal-proportions method since 1941
(2~U.S.C.~\usca{}~2a(a)).\footnote{%
  The method of equal proportions has been used for every apportionment since
  1941.
  The Supreme Court upheld it in \textit{United States Department of Commerce
  v.\ Montana}, 503 U.S. 442 (1992), noting that while the Constitution does
  not mandate a specific mathematical method, equal proportions ``minimizes
  the percentage differences between the average district populations.''
}
No one disputes that Huntington-Hill is the correct method for the interstate
problem.
It is law.

\paragraph{Step~3: Politicians should not draw their own districts.}
Both parties have articulated this principle when the other party holds the pen.
Republicans in Maryland, Democrats in North Carolina, Republicans in New York---
each has argued, on the record, in court, that self-interested map-drawing
is antidemocratic.\footnote{%
  See \rucho{}, 588 U.S. at 689--690 (cataloguing bipartisan complaints);
  \textit{Common Cause v.\ Lewis}, N.C. Super. Ct., 2019 (Democratic plaintiffs
  challenging Republican-drawn maps using bipartisan reform arguments).
}
Independent redistricting commissions have been adopted in seventeen states
for exactly this reason.
This is not a partisan point.
It is a point about institutional integrity.

\paragraph{Step~4: If Huntington-Hill is right for the interstate problem,
it is right for the structurally identical intrastate problem.}
The interstate problem is: given a state's population, allocate
$k$~congressional seats to the state in proportion to population.
The intrastate problem is: given a state's geographic units (census tracts),
allocate $k$~congressional districts to those units in proportion to population.
Both problems require apportioning $k$ indivisible units among sub-units
in proportion to a weight (population).
Both require minimizing the percentage deviation between actual and ideal
allocation.
The problems are isomorphic.

If \hh{} is correct for the interstate problem---and it is, because Congress
said so and no one has challenged it in eighty years---then \hh{} is correct
for the intrastate problem by the same logic.
No new constitutional theory is required.
No new doctrine needs to be invented.
The same principle, applied at the next level down.

\paragraph{The conclusion.}
Apply \hh{} within states via \ar{}.
The prime factorization of the seat count determines the bisection tree.
The TIGER adjacency graph determines the boundaries.
No political choices remain.
The \dia{} encodes this in one sentence.

\subsection{Scope of This Article}

This Article proceeds as follows.
Section~\ref{sec:first-principles} develops the four-step argument in full,
addressing each premise and the logical connection between them.
Section~\ref{sec:why-apportionregions} explains why \ar{} is the uniquely
correct intrastate implementation of \hh{}---and why the other algorithms
in the B-series toolbox, however sophisticated, are practitioner choices
rather than constitutionally derived ones.
Section~\ref{sec:bipartisan} addresses, separately, the specific objections
that Republicans and Democrats have historically raised against federal
redistricting mandates, and shows why the \dia{} satisfies each party's
legitimate interests.
Section~\ref{sec:implementation} details the statute's implementation:
the \hh{} bisection tree, the VRA overlay for Gingles-qualified states,
the county-preservation default, the seed formula, and the fallback for
states where the seat count's factorization is incompatible with a
strict \hh{} tree.
Section~\ref{sec:conclusion} restates the argument, presents the proposed
statutory text, and explains why this is the right moment---after \callais{},
after \rucho{}, after a decade of state-court redistricting chaos---to enact it.

A word on tone.
This Article is an advocacy piece, not a survey.
It argues a specific position: the \dia{} should pass, in the form
proposed in Section~\ref{sec:conclusion}.
The Article presents objections fairly and responds to them.
But it does not pretend to be agnostic.
The argument from first principles is, in the author's view, airtight.
The reader is invited to find the flaw.
